% ============================================================
%  PRELIMINARES: DEDICATORIA, AGRADECIMIENTOS Y RESUMEN
%  ------------------------------------------------------------
%  El glosario vive en secciones/11_glosario.tex (después de los
%  índices, orden §2.2) y el índice general, el de tablas y el de
%  figuras se generan desde main.tex. El orden completo de
%  Preliminares lo decide main.tex, no este archivo.
% ============================================================

% ---------- DEDICATORIA -------------------------------------
% No muestra número de página, pero SÍ se cuenta (regla 3.1).
\thispagestyle{empty}
\vspace*{6cm}
\begin{flushright}
  \interlineadoResumen
  \textit{A mi abuela, a toda mi familia y amigos que me enseñaron a no darme
   por vencido y el valor de la paciencia y la persistencia;  
   a Fer por el apoyo académico que me brindó durante casi toda la carrera,
   a las comunidades que me dieron el criterio para poder crecer por mi cuenta.}
\end{flushright}
\clearpage

% ---------- AGRADECIMIENTOS ---------------------------------
\chapter*{AGRADECIMIENTOS}
\thispagestyle{empty}   % debe ir DESPUÉS de \chapter*, que reimpone el estilo
\begingroup           % \interlineadoResumen no debe filtrarse al resto del documento
\interlineadoResumen
Agradezco a mi familia, que me sostiene, impulsa y ayuda sin importar las adversidades.
Agradezco también a la vida por todas aquellas personas que por azares del destino llegué
a conocer, con las cuales he pasado inolvidables momentos. Agradezco a los docentes del
Diplomado en Ciencia de Datos, a Ismael y al equipo de Trufi Association por facilitar los datos que
hicieron posible este trabajo.
\endgroup
\clearpage

% ---------- RESUMEN -----------------------------------------
% Formato exigido: Garamond 11.5 pt, interlineado múltiple 1.2,
% con espaciado anterior y posterior de 6 puntos entre párrafos.
% Extensión: idealmente 1 plana COMPLETA incluyendo las palabras
% clave. Estructura sugerida por la guía:
%   1 párrafo de introducción
%   1-2 párrafos de metodología
%   1 párrafo de hallazgos (resultados + breve análisis)
%   1 párrafo de conclusiones
\chapter*{RESUMEN}
\addcontentsline{toc}{chapter}{RESUMEN}
{%
  \garamondfont\fontsize{11.5}{13.8}\selectfont   % Garamond 11.5, interlineado 1.2
  \setlength{\parskip}{6pt}                       % 6 pt antes/después
  \setlength{\parindent}{0pt}

  % --- Párrafo 1: introducción / contexto ---------------------
  El transporte público del eje metropolitano de Cochabamba opera principalmente bajo la
  modalidad de paratránsito, y Trufi App constituye una de las pocas fuentes públicas de
  información sobre sus rutas, mantenida mediante el trabajo voluntario de Trufi
  Association. La aplicación registra cuántas consultas de ruta se originan en cada zona,
  pero no cuántas cabría esperar según sus características; en consecuencia, un conteo
  bajo no permite distinguir entre una zona poco poblada y una zona donde la aplicación no
  resulta útil. La presente monografía estima el número esperado de consultas por celda
  hexagonal H3 de resolución 8 (aproximadamente 0,74~km\textsuperscript{2}) a partir de la
  población, la ubicación y el contexto territorial de cada celda, utilizando 1.927.675
  consultas registradas entre septiembre de 2022 y junio de 2024.

  % --- Párrafo 2: metodología ----------------------------------
  El proyecto se desarrolló bajo la metodología CRISP-DM. Tras la depuración de los datos
  se conservaron 1.924.578 consultas válidas y se construyó una tabla de 1.381 celdas
  pobladas según Kontur Population, que incluye 444 celdas sin ninguna consulta
  registrada. Dado que el análisis exploratorio evidenció una fuerte autocorrelación
  espacial de la demanda, antes de ajustar cualquier modelo se reservó como conjunto de
  prueba el 20~\% de los bloques espaciales. Se compararon siete técnicas mediante
  validación cruzada espacial por bloques: cuatro líneas base de tasa (global, por
  municipio, por anillo de distancia y por vecindad H3), dos modelos lineales generalizados
  con offset poblacional (Poisson y Binomial Negativa) y un modelo de
  gradient boosting con pérdida Poisson. La técnica final se seleccionó mediante
  una regla declarada antes de conocer los resultados.

  % --- Párrafo 3: hallazgos ------------------------------------
  La regla adoptó la estimación por vecindad H3, que multiplica la tasa de consultas por
  habitante de las celdas vecinas por la población de cada celda. En el conjunto de
  prueba, evaluado una única vez, esta técnica redujo la devianza de Poisson de 1.807,2 a
  981,7 respecto de una tasa única para toda el área y explicó el 76,5~\% de la devianza.
  Los modelos basados en población y distancia al centro no la superaron, y una validación
  aleatoria habría sobrestimado la capacidad predictiva (D\textsuperscript{2} de 0,91
  frente a 0,64). La brecha entre consultas observadas y esperadas resultó estable ante
  variaciones razonables de los datos (correlación de Spearman $\geq 0{,}87$), y las
  celdas sin una parada de la red mapeada a 500~m o menos presentaron, de forma
  sistemática, menos consultas de las esperadas ($\delta$ de Cliff $= +0{,}37$).

  % --- Párrafo 4: conclusiones ---------------------------------
  Se concluye que es posible estimar la demanda esperada de información de transporte por
  celda con un error, en zonas no vistas, cercano a la mitad del de una tasa uniforme, y
  que la mejor estimación disponible es de carácter local. El resultado se entrega a Trufi
  Association como un mapa interactivo, un archivo de predicciones por celda y una lista de
  veinte celdas prioritarias para revisión en terreno, concebidos como apoyo a la
  priorización del mapeo y no como evidencia concluyente de ausencia de rutas.

  \vspace{6pt}
  \noindent\textbf{Palabras clave:} demanda esperada, datos de conteo, validación cruzada
  espacial, paratránsito, Trufi App, H3, Kontur Population, CRISP-DM.
}
\clearpage
